\section{Legge di Bayes}

Consideriamo due variabili aleatorie $X$ e $Y$ con PMF congiunta $p_{X,Y}(x,y)$.

Consideriamo l'evento:
\[
\{X=x, Y=y\} = \{X=x\} \cap \{Y=y\}.
\]

Allora:
\begin{align*}
P(X=x, Y=y)
&= P(\{X=x\} \cap \{Y=y\}) \\
&= P(X=x \mid Y=y)\cdot P(Y=y) \\
&= P(Y=y \mid X=x)\cdot P(X=x).
\end{align*}

Scrivendo in termini di PMF:
\[
p_{X|Y}(x|y)\cdot p_Y(y) = p_{Y|X}(y|x)\cdot p_X(x).
\]

Da cui otteniamo la \textbf{Legge di Bayes}:
\begin{equation}
\boxed{
p_{X|Y}(x|y) = \frac{p_{Y|X}(y|x)\cdot p_X(x)}{p_Y(y)}
}
\end{equation}

\subsection{Esempio: Test diagnostico}

Immaginiamo la seguente situazione.

Svegli e ti senti malato, vai dal dottore e non sai se sei malato oppure no. Il dottore ti fa fare un test.

Supponiamo che il test rilevi correttamente la malattia nel $99\%$ dei casi, e sbagli solo nell'$1\%$ dei casi.

La malattia colpisce solo lo $0.1\%$ della popolazione.

Il dottore chiede:
\textit{``Qual è la probabilità che tu abbia la malattia se il test è risultato positivo?''}

\subsubsection{Formalizzazione}

Definiamo:

\begin{itemize}
\item $H$: evento \textit{``la persona ha la malattia''}.
\item $E$: evento \textit{``test positivo''}.
\end{itemize}

Applichiamo Bayes:
\begin{equation}
P(H|E) = \frac{P(E|H)\cdot P(H)}{P(E)}.
\end{equation}

Dove:
\begin{itemize}
\item $P(E|H)$: probabilità che il test sia positivo se la persona è malata.
\item $P(H)$: probabilità che una persona sia malata.
\item $P(E)$: probabilità che il test risulti positivo.
\end{itemize}

Ora calcoliamo $P(E)$ usando la legge della probabilità totale:
\begin{equation}
P(E) = P(H)\cdot P(E|H) + P(\overline{H})\cdot P(E|\overline{H}).
\end{equation}

Sostituendo:
\[
P(H) = 0.001, \quad
P(E|H) = 0.99, \quad
P(E|\overline{H}) = 0.01.
\]

Otteniamo:
\[
P(H|E) = \frac{0.99\cdot 0.001}{0.99\cdot 0.001 + 0.999\cdot 0.01}
= \frac{0.00099}{0.00099 + 0.00999}
\approx 0.09 = 9\%.
\]

\subsubsection{Interpretazione Intuitiva}

Supponiamo di avere $1000$ persone.

\begin{itemize}
\item Una sola persona ha la malattia.
\item Tra le $999$ persone sane, l'$1\%$ risulta comunque positivo: circa $10$ persone.
\end{itemize}

Quindi abbiamo:

\begin{itemize}
\item $1$ vero positivo,
\item $10$ falsi positivi.
\end{itemize}

In totale, tra i positivi, ci sono $11$ persone, di cui solo una è davvero malata.

\[
P(\text{malato} \mid \text{positivo}) = \frac{1}{11} \approx 9\%.
\]

\subsubsection{Ripetizione del Test}

Cosa succede se ripetiamo il test in un altro laboratorio?

Supponiamo che i due test siano indipendenti.

La probabilità aggiornata diventa nuovamente:
\[
P(H|E) = \frac{P(E|H)\cdot P(H)}{P(H)\cdot P(E|H) + P(\overline{H})\cdot P(E|\overline{H})}.
\]

Dopo il secondo test positivo, la probabilità di essere davvero malati sale fino a circa:

\[
\boxed{91\%}.
\]

\subsection{Osservazione Importante}

Questo esempio mostra come la probabilità iniziale (\textit{prior}) influenzi fortemente il risultato finale.

Un test anche molto accurato può produrre una probabilità finale sorprendentemente bassa se la malattia è molto rara.

\begin{quote}
\textit{``Un test molto preciso non implica automaticamente una diagnosi affidabile.''}
\end{quote}